% ID: SEL-D0038-OB-E03
% ESTADO: cotejada. Segunda pasada visual y borrador compilado revisados.
% Fuente: Selectividad/matematicas_aplicadas/Ua3MATEMÁTICAS APLICADAS examen.pdf
% SHA-256: 4440c15a31a77ee9e052a0668923508e117795b369d569bc02d24b7937366f7a
% Página: 2; opción B; ejercicio original 3.

En una urna A hay 8 bolas verdes y 6 rojas. En otra urna B hay 4 bolas verdes, 5 rojas y 1 negra. Se lanza un dado, si sale un número menor que 3 se saca una bola de la urna A, y si sale mayor o igual que 3 se saca una bola de la urna B.
\begin{enumerate}[label=\alph*)]
\item \textbf{(0.5 puntos)} Calcule la probabilidad de que la bola sea verde si ha salido un 4.
\item \textbf{(1 punto)} Calcule la probabilidad de que la bola elegida sea roja.
\item \textbf{(1 punto)} Sabiendo que ha salido una bola verde, \textquestiondown{}cuál es la probabilidad de que sea de la urna A?
\end{enumerate}
